\documentclass[11pt,a4paper]{article}

\usepackage[utf8]{inputenc}
\usepackage[T1]{fontenc}
\usepackage[french]{babel}
\usepackage{amsmath,amssymb,amsfonts}
\usepackage{booktabs}
\usepackage{geometry}
\usepackage{graphicx}
\usepackage{hyperref}
\usepackage{float}
\usepackage{caption}
\usepackage{subcaption}

\geometry{margin=2.4cm}
\hypersetup{colorlinks=true, linkcolor=blue, urlcolor=blue, citecolor=blue}

\title{Deep Hedging d'un Put européen avec perte quadratique et prime connue}
\author{Projet académique de niveau master/ingénieur}
\date{\today}

\begin{document}
\maketitle

\begin{abstract}
Ce rapport étudie la couverture d'un Put européen par Deep Hedging dans un
marché discrétisé. La prime est supposée connue et l'objectif est de minimiser
une perte quadratique du P\&L terminal, conformément à la formulation à prime
exogène de Buehler, Gonon, Teichmann et Wood. Les stratégies sont paramétrées
par réseaux de neurones et entraînées par rétropropagation à travers des
trajectoires Monte Carlo. Les performances sont comparées à un delta hedging
Black--Scholes et analysées selon l'architecture, les paramètres de marché et
la présence de coûts proportionnels.
\end{abstract}

\section{Introduction}

Le sujet imposé est le suivant : implémenter un algorithme de Deep Hedging
inspiré de l'article fourni pour couvrir un Put, en supposant une loss
quadratique et une prime connue, puis analyser la qualité des résultats en
fonction de l'architecture du réseau. Le livrable doit être un dépôt Python
reproductible contenant code modulaire, notebook exécutable, rapport et
figures.

L'article de Buehler et al. propose de remplacer la recherche analytique d'une
stratégie de couverture par un problème d'optimisation statistique sur des
stratégies paramétrées. Son point important pour ce projet est la distinction
entre deux usages : l'indifference pricing sous mesures de risque convexes et
le cas, explicitement mentionné dans leur remarque 3.4, où la prime $p_0$ est
donnée par le marché. Nous retenons ce second cadre : la prime n'est pas
optimisée, et la stratégie de hedge minimise directement une perte terminale.

Le choix d'un Put européen est volontairement simple : il permet de disposer
d'un benchmark Black--Scholes clair, tout en conservant les enjeux de couverture
discrète, de simulation sous probabilité réelle et de comparaison
hors-échantillon. Le projet ne cherche pas à paraphraser l'article ; il adapte
sa structure de contrôle stochastique au cas pédagogique d'une option vanille
avec perte quadratique.

\section{Cadre mathématique}

\subsection{Espace probabilisé, filtration et dates de décision}

On travaille en temps discret sur une grille de couverture
\[
0=t_0<t_1<\cdots<t_n=T,\qquad t_k=k\Delta t,\qquad \Delta t=\frac{T}{n}.
\]
Le modèle est défini sur un espace probabilisé filtré
\[
(\Omega,\mathcal{F},(\mathcal{F}_{t_k})_{k=0,\ldots,n},\mathbb{P}),
\]
où $\mathcal{F}_{t_k}$ représente toute l'information observable au moment où
la décision de couverture $\delta_k$ est prise. La mesure $\mathbb{P}$ est la
probabilité réelle, ou historique : elle décrit la distribution statistique des
scénarios qui servent à entraîner et à évaluer la couverture. Cette distinction
est importante car le problème traité ici n'est pas un problème de pricing pur,
mais un problème de minimisation d'erreur de couverture sous la loi de
réalisation des marchés.

Dans l'implémentation, l'information disponible à $t_k$ est résumée par des
variables d'état observables :
\[
I_k=\left(S_{t_k},\, t_k,\, K,\, T,\, \delta_{k-1}\right),
\]
ou, après normalisation numérique,
\[
X_k=\left(\frac{S_{t_k}}{S_0}-1,\,
\log\frac{S_{t_k}}{K},\,
\frac{T-t_k}{T},\,
\delta_{k-1}\right).
\]
Le choix de $X_k$ est cohérent avec un Put européen dans un modèle markovien :
le prix courant, la moneyness et le temps restant suffisent à représenter l'état
de marché pertinent, tandis que $\delta_{k-1}$ est nécessaire dès que l'on
introduit des coûts de transaction ou des contraintes de turnover.

\subsection{Dynamique du sous-jacent sous la probabilité réelle}

Le sous-jacent $S$ est simulé sous $\mathbb{P}$ par un modèle de Black--Scholes
discrétisé exactement. Pour $k=0,\ldots,n-1$,
\[
S_{t_{k+1}}
=S_{t_k}\exp\left[
\left(\mu-\frac12\sigma^2\right)\Delta t
+\sigma\sqrt{\Delta t}\,\varepsilon_{k+1}
\right],
\]
avec
\[
\varepsilon_{1},\ldots,\varepsilon_n
\stackrel{i.i.d.}{\sim}\mathcal{N}(0,1)
\quad\text{sous }\mathbb{P}.
\]
Ainsi
\[
\log\frac{S_{t_{k+1}}}{S_{t_k}}
\sim
\mathcal{N}\left(
\left(\mu-\frac12\sigma^2\right)\Delta t,\,
\sigma^2\Delta t
\right).
\]
Le paramètre $\mu$ est la dérive réelle utilisée pour la génération de scénarios.
Il n'intervient pas dans le prix Black--Scholes de référence, mais il influence
la distribution des erreurs de couverture hors-échantillon lorsque l'objectif
est évalué sous $\mathbb{P}$.

\subsection{Mesure risque-neutre et rôle limité du pricing}

Pour fixer une prime de référence et construire le benchmark de delta hedging,
on introduit séparément une mesure risque-neutre $\mathbb{Q}$ sous laquelle le
processus actualisé est martingale. Dans le cas Black--Scholes avec taux sans
risque constant $r$,
\[
S_{t_{k+1}}
=S_{t_k}\exp\left[
\left(r-\frac12\sigma^2\right)\Delta t
+\sigma\sqrt{\Delta t}\,\varepsilon^{\mathbb{Q}}_{k+1}
\right],
\qquad
\varepsilon^{\mathbb{Q}}_{k+1}\sim\mathcal{N}(0,1).
\]
La prime connue $p_0$ peut être imposée extérieurement. Par défaut, elle est
prise égale au prix risque-neutre Black--Scholes du Put :
\[
p_0=P_{\mathrm{BS}}(S_0,K,T,r,\sigma)
=Ke^{-rT}N(-d_2)-S_0N(-d_1),
\]
où
\[
d_1=\frac{\log(S_0/K)+(r+\frac12\sigma^2)T}{\sigma\sqrt{T}},
\qquad
d_2=d_1-\sigma\sqrt{T}.
\]
Le delta du Put utilisé comme benchmark est
\[
\Delta_{\mathrm{put}}(t,S_t)
=\frac{\partial P_{\mathrm{BS}}}{\partial S}(t,S_t)
=N(d_1(t,S_t))-1,
\]
avec
\[
d_1(t,S)
=\frac{\log(S/K)+(r+\frac12\sigma^2)(T-t)}
{\sigma\sqrt{T-t}}.
\]
Cette partie relève de $\mathbb{Q}$. En revanche, l'entraînement Deep Hedging
minimise une espérance sous $\mathbb{P}$. La méthode ne suppose donc pas que
la stratégie optimale soit la delta Black--Scholes ; cette dernière est seulement
un point de comparaison naturel en marché complet sans frictions et en temps
continu.

\subsection{Payoff, convention vendeur et stratégie admissible}

Le produit couvert est un Put européen de strike $K$ et de maturité $T$ :
\[
Z=g(S_T)=(K-S_T)^+.
\]
On adopte la convention du vendeur de l'option : à la date initiale, il reçoit
la prime connue $p_0$ ; à maturité, il verse $Z$ ; entre les deux, il peut
échanger le sous-jacent. Une stratégie de couverture est un processus
\[
\delta=(\delta_0,\ldots,\delta_{n-1})
\]
tel que $\delta_k$ désigne le nombre d'unités du sous-jacent détenues sur
l'intervalle $[t_k,t_{k+1})$. La contrainte fondamentale est la prévisibilité :
\[
\delta_k\ \text{est}\ \mathcal{F}_{t_k}\text{-mesurable}.
\]
Autrement dit, $\delta_k$ ne peut pas dépendre de $S_{t_{k+1}}$ ni d'aucune
information future. L'ensemble des stratégies admissibles est noté
$\mathcal{H}$. Dans la version de base, $\mathcal{H}$ contient les stratégies
prévisibles bornées générées par les réseaux, la borne provenant de la sortie
en $\tanh$. Cette borne joue le rôle d'une contrainte de levier numérique :
\[
|\delta_k|\leq \bar{\delta}.
\]

\subsection{Autofinancement, gain de trading et coûts}

En absence de frictions, le gain cumulé d'une stratégie prévisible est
l'intégrale stochastique discrète
\[
(\delta\cdot S)_T
=\sum_{k=0}^{n-1}\delta_k\left(S_{t_{k+1}}-S_{t_k}\right).
\]
Cette écriture encode l'autofinancement : les rééquilibrages intermédiaires
sont financés par le portefeuille lui-même, et seul le P\&L terminal est évalué.

Pour l'extension avec frictions, on retient un coût proportionnel au montant
notionnel échangé. En posant
\[
\delta_{-1}=0,\qquad \delta_n=0,
\]
ce qui signifie absence de position avant l'ouverture et liquidation finale à
maturité, le coût total est
\[
C_T(\delta)
=\lambda\sum_{k=0}^{n}S_{t_k}\left|\delta_k-\delta_{k-1}\right|.
\]
Le terme $k=0$ correspond à l'ouverture de la position, les termes
$1,\ldots,n-1$ aux rééquilibrages, et le terme $k=n$ à la liquidation. Dans
le cas principal sans coûts, $\lambda=0$ et donc $C_T(\delta)=0$.

\subsection{P\&L terminal et erreur de couverture}

Le P\&L terminal du vendeur couvert est
\[
\Pi_T(\delta)
=-Z+p_0+(\delta\cdot S)_T-C_T(\delta).
\]
Une couverture parfaite en temps discret satisferait $\Pi_T(\delta)=0$ presque
sûrement. En pratique, cette égalité n'est pas atteinte pour trois raisons :
\begin{itemize}
\item les dates de couverture sont discrètes ;
\item la stratégie est restreinte à une famille paramétrique ;
\item les coûts de transaction ou contraintes de position rendent le marché
effectivement incomplet.
\end{itemize}
L'erreur de couverture est donc identifiée au P\&L résiduel $\Pi_T$. Un P\&L
positif signifie que le vendeur termine avec un excédent après paiement du Put ;
un P\&L négatif signifie une perte résiduelle.

\subsection{Objectif quadratique avec prime connue}

La prime $p_0$ n'est pas optimisée : elle est considérée comme donnée par le
marché ou par une convention de pricing. Le problème est alors celui de la
meilleure stratégie de réduction du risque résiduel :
\[
\delta^\star
\in
\operatorname*{arg\,min}_{\delta\in\mathcal{H}}
J(\delta),
\qquad
J(\delta)
=
\mathbb{E}_{\mathbb{P}}
\left[
\ell\left(\Pi_T(\delta)\right)
\right],
\]
avec la perte quadratique
\[
\ell(x)=x^2.
\]
Explicitement,
\[
J(\delta)
=
\mathbb{E}_{\mathbb{P}}
\left[
\left(
-(K-S_T)^+
+p_0
+\sum_{k=0}^{n-1}\delta_k(S_{t_{k+1}}-S_{t_k})
-C_T(\delta)
\right)^2
\right].
\]
Cette fonction objectif correspond à une erreur quadratique moyenne de
réplication. Elle pénalise symétriquement pertes et gains résiduels : l'objectif
est de centrer et concentrer la distribution du P\&L autour de zéro, non de
maximiser le gain espéré du vendeur. Cette propriété est adaptée à une étude
académique de qualité de couverture. Elle diffère d'une mesure prudentielle
asymétrique, telle que la CVaR, qui pénalise spécifiquement la queue gauche du
P\&L.

Lorsque $\lambda=0$ et que le modèle Black--Scholes est utilisé en temps
continu avec trading continu, la couverture delta réplique l'option. Ici, même
avec le même modèle de diffusion, le passage au temps discret laisse une erreur
de discrétisation. Le Deep Hedging apprend directement la stratégie qui réduit
cette erreur sous la loi simulée et sous les contraintes numériques choisies.

\section{Méthodologie Deep Hedging}

\subsection{Point de départ dans l'article de Buehler et al.}

Buehler, Gonon, Teichmann et Wood formulent le Deep Hedging comme un problème
d'optimisation de stratégies adaptées dans un marché discret avec frictions
éventuelles. Leur cadre général considère une position terminale
\[
-Z+p_0+(\delta\cdot S)_T-C_T(\delta)
\]
et remplace la recherche analytique d'une couverture par l'optimisation
numérique d'une stratégie paramétrée. Deux idées de l'article sont reprises ici.

Premièrement, la stratégie est vue comme un contrôle prévisible : à chaque date,
une fonction de décision transforme l'information disponible en position de
portefeuille. Deuxièmement, l'objectif est évalué directement sur la distribution
des scénarios simulés, ce qui permet de traiter les frictions sans dériver de
PDE ou de formule fermée.

Le point exact utilisé pour ce projet est le cas à prime exogène. Si $p_0$ est
donné, l'article conduit à l'optimisation
\[
\inf_{\delta\in\mathcal{H}}
\mathbb{E}_{\mathbb{P}}
\left[
\ell\left(
-Z+p_0+(\delta\cdot S)_T-C_T(\delta)
\right)
\right].
\]
Dans notre étude,
\[
\ell(x)=x^2,
\]
ce qui transforme le Deep Hedging en minimisation directe de la MSE de
couverture. La partie prix d'indifférence de l'article n'est donc pas l'objet
principal : on ne cherche pas $p_0$, on cherche $\delta^\star$ conditionnellement
à $p_0$.

\subsection{Formulation comme contrôle stochastique discret}

Le problème peut se lire comme un problème de contrôle stochastique à horizon
fini. L'état contrôlé est constitué du marché observable et de la position
précédente :
\[
X_k=\Phi(S_{t_k},t_k,K,T,\delta_{k-1}),
\]
où $\Phi$ désigne la normalisation utilisée numériquement. L'action choisie à
la date $t_k$ est
\[
a_k=\delta_k.
\]
La dynamique contrôlée du portefeuille ne modifie pas la dynamique de $S$ dans
la version de base, car on ne modélise pas d'impact de marché. Le contrôle
agit seulement sur le gain de trading et sur les coûts :
\[
W_{k+1}
=W_k+\delta_k(S_{t_{k+1}}-S_{t_k})
-c_k(\delta_k-\delta_{k-1}),
\]
avec $W_0=p_0$ et, à maturité,
\[
\Pi_T=W_T-Z-c_n(0-\delta_{n-1}).
\]
Pour le coût proportionnel,
\[
c_k(u)=\lambda S_{t_k}|u|.
\]
En absence de coûts, $c_k\equiv 0$.

Une formulation dynamique exacte ferait intervenir des fonctions valeur
\[
V_k(x,w,\delta_{k-1})
=
\inf_{\delta_k,\ldots,\delta_{n-1}}
\mathbb{E}_{\mathbb{P}}
\left[
\Pi_T^2
\mid X_k=x,\ W_k=w,\ \delta_{k-1}
\right].
\]
Cependant, même dans un modèle simple, cette récursion devient coûteuse dès
que l'état est enrichi ou que des frictions sont ajoutées. Le Deep Hedging évite
de résoudre explicitement cette programmation dynamique : il optimise
directement une politique globale paramétrée sur des trajectoires Monte Carlo.

\subsection{Paramétrisation neuronale des stratégies}

La stratégie admissible est restreinte à une famille paramétrique
\[
\mathcal{H}_{\Theta}
=
\left\{
\delta^\theta:\theta\in\Theta
\right\}
\subset\mathcal{H}.
\]
Pour un MLP partagé entre les dates, la décision prend la forme
\[
\delta_k^\theta
=
\bar{\delta}\,
\tanh\left(F_\theta(X_k)\right),
\]
où $F_\theta$ est un réseau feed-forward et $\bar{\delta}>0$ borne la position.
L'usage du même réseau à toutes les dates impose une forme de régularisation :
le temps restant est fourni en entrée, ce qui permet au réseau d'apprendre une
politique dépendante de $k$ sans multiplier les paramètres par le nombre de
dates.

Pour un réseau à $L$ couches, une écriture typique est
\[
h_0=X_k,\qquad
h_{\ell}=\rho(A_{\ell}h_{\ell-1}+b_{\ell}),
\quad \ell=1,\ldots,L-1,
\]
\[
F_\theta(X_k)=A_Lh_{L-1}+b_L,
\]
avec activation ReLU $\rho(x)=\max(x,0)$ dans l'implémentation. Les paramètres
sont
\[
\theta=\{A_\ell,b_\ell\}_{\ell=1,\ldots,L}.
\]
Le MLP simple possède une faible capacité et sert de référence neuronale. Le
MLP profond augmente le nombre de couches et la dimension cachée pour mieux
approximer une fonction de couverture non linéaire en spot, moneyness et temps.

Pour le LSTM, on fournit la séquence
\[
(Y_0,\ldots,Y_k),
\qquad
Y_j=\left(\frac{S_{t_j}}{S_0}-1,\log\frac{S_{t_j}}{K},\frac{T-t_j}{T}\right),
\]
à un état récurrent
\[
H_k=\mathrm{LSTM}_\theta(Y_0,\ldots,Y_k).
\]
La position est ensuite produite par une tête feed-forward utilisant $H_k$ et
$\delta_{k-1}$ :
\[
\delta_k^\theta
=
\bar{\delta}\,
\tanh\left(G_\theta(H_k,\delta_{k-1})\right).
\]
Dans un Black--Scholes markovien, le LSTM n'est pas théoriquement nécessaire ;
son intérêt est expérimental : tester si une mémoire séquentielle stabilise
l'apprentissage ou capte des effets de chemin lorsque l'on enrichit le modèle.

\subsection{Objectif empirique Monte Carlo}

L'espérance $J(\theta)$ n'est pas calculée analytiquement. On génère
$M$ trajectoires indépendantes
\[
\left(S^{(m)}_{t_0},\ldots,S^{(m)}_{t_n}\right),
\qquad m=1,\ldots,M,
\]
et l'on remplace $J$ par l'estimateur empirique
\[
\widehat{J}_M(\theta)
=
\frac1M
\sum_{m=1}^{M}
\left[
\Pi_T^{(m)}(\theta)
\right]^2,
\]
où
\[
\Pi_T^{(m)}(\theta)
=
-\left(K-S_T^{(m)}\right)^+
+p_0
+\sum_{k=0}^{n-1}
\delta_k^\theta\!\left(X_k^{(m)}\right)
\left(S_{t_{k+1}}^{(m)}-S_{t_k}^{(m)}\right)
-C_T^{(m)}(\delta^\theta).
\]
En pratique, l'optimisation utilise des mini-batches $B\subset\{1,\ldots,M\}$ :
\[
\widehat{J}_B(\theta)
=
\frac1{|B|}
\sum_{m\in B}
\left[
\Pi_T^{(m)}(\theta)
\right]^2.
\]
Cette approximation introduit un bruit de gradient, mais elle rend
l'optimisation scalable et joue un rôle régularisant usuel en apprentissage
profond.

\subsection{Backpropagation à travers les trajectoires}

Pour une trajectoire donnée, le P\&L est une composition de fonctions
différentiables des paramètres $\theta$ :
\[
\theta
\mapsto
\delta_0^\theta,\ldots,\delta_{n-1}^\theta
\mapsto
\Pi_T(\theta)
\mapsto
\Pi_T(\theta)^2.
\]
Le gradient empirique s'écrit formellement
\[
\nabla_\theta \widehat{J}_B(\theta)
=
\frac{2}{|B|}
\sum_{m\in B}
\Pi_T^{(m)}(\theta)
\nabla_\theta \Pi_T^{(m)}(\theta).
\]
Sans coûts de transaction,
\[
\nabla_\theta \Pi_T^{(m)}(\theta)
=
\sum_{k=0}^{n-1}
\left(S_{t_{k+1}}^{(m)}-S_{t_k}^{(m)}\right)
\nabla_\theta \delta_k^\theta(X_k^{(m)}).
\]
Avec coûts proportionnels, des termes supplémentaires apparaissent :
\[
-\lambda
\sum_{k=0}^{n}
S_{t_k}^{(m)}
\nabla_\theta
\left|
\delta_k^\theta-\delta_{k-1}^\theta
\right|.
\]
La valeur absolue n'est pas différentiable en zéro, mais les bibliothèques
d'autodifférentiation utilisent un sous-gradient, ce qui est suffisant pour
l'optimisation numérique. PyTorch construit automatiquement le graphe de calcul
sur toutes les dates de couverture et applique la règle de la chaîne.

La mise à jour des paramètres est effectuée par Adam :
\[
\theta_{j+1}
=
\theta_j-\eta_j\,\widehat{m}_j/(\sqrt{\widehat{v}_j}+\epsilon),
\]
où $\widehat{m}_j$ et $\widehat{v}_j$ sont les moments adaptatifs du gradient.
Un clipping de gradient est ajouté pour éviter des mises à jour extrêmes au
début de l'entraînement.

\subsection{Lien avec le delta hedging classique}

Le delta hedging Black--Scholes correspond à la stratégie analytique
\[
\delta_k^{\mathrm{BS}}
=
\Delta_{\mathrm{put}}(t_k,S_{t_k}).
\]
En marché complet, sans coûts et avec couverture continue, cette stratégie
réplique le Put. En temps discret, elle devient un benchmark non parfait :
\[
\Pi_T(\delta^{\mathrm{BS}})
\neq 0
\quad\text{en général}.
\]
Le Deep Hedging ne cherche pas à apprendre la formule du delta pour elle-même ;
il minimise l'erreur quadratique terminale dans le dispositif discret réellement
utilisé. Il peut donc, en principe, s'écarter du delta si cela réduit la MSE
sous $\mathbb{P}$, avec coûts, contraintes ou architecture donnée.

\subsection{Architectures comparées}

Trois familles de politiques sont comparées expérimentalement :
\begin{itemize}
\item \textbf{MLP simple} : une couche cachée. Il permet de tester si une faible
capacité suffit à approximer une couverture proche du delta dans un cadre
markovien simple.
\item \textbf{MLP profond} : plusieurs couches cachées. Il augmente la capacité
d'approximation et doit mieux représenter les non-linéarités liées à la
moneyness, au temps restant et aux coûts.
\item \textbf{LSTM} : politique récurrente. Elle est surtout utile lorsque le
modèle ou les features ne sont plus strictement markoviens ; dans le cas
Black--Scholes, elle sert de test de robustesse architectural.
\end{itemize}

\section{Implémentation}

Le code est organisé en quatre modules :
\begin{itemize}
\item \texttt{simulation.py} : paramètres de marché, simulation GBM,
payoff, prix et delta Black--Scholes.
\item \texttt{model.py} : politiques MLP et LSTM, fabrique de modèles.
\item \texttt{training.py} : P\&L, coûts, fonctions de perte, entraînement,
évaluation et benchmark delta.
\item \texttt{utils.py} : tableaux et figures.
\end{itemize}

Les trajectoires Monte Carlo utilisent l'échantillonnage antithétique :
chaque choc gaussien $\varepsilon$ est apparié à $-\varepsilon$. Cette méthode
ne change pas la loi marginale mais réduit la variance des estimations pour un
coût de calcul inchangé. Les seeds NumPy et PyTorch sont fixées pour assurer la
reproductibilité.

La prime par défaut est
\[
p_0=K e^{-rT}N(-d_2)-S_0N(-d_1),
\]
avec $d_1,d_2$ standards de Black--Scholes sous $\mathbb{Q}$. Une prime
externe peut être imposée en remplaçant ce paramètre.

\section{Résultats numériques}

Les figures et tableaux du notebook sont générés dans \texttt{report/figures}.
Le tableau principal compare les stratégies selon la moyenne du P\&L, la MSE,
la RMSE, la variance et un intervalle de confiance normal à 95\% pour la MSE.
Un exemple de structure attendue est donné ci-dessous ; les valeurs exactes
dépendent du nombre d'époques choisi lors de l'exécution.

\begin{table}[H]
\centering
\caption{Mesures hors-échantillon à reporter après exécution du notebook.}
\begin{tabular}{lrrrr}
\toprule
Stratégie & Moyenne P\&L & MSE & RMSE & Variance\\
\midrule
Delta Black--Scholes & -- & -- & -- & --\\
MLP simple & -- & -- & -- & --\\
MLP profond & -- & -- & -- & --\\
LSTM & -- & -- & -- & --\\
\bottomrule
\end{tabular}
\end{table}

Dans un marché Black--Scholes sans coûts, le delta hedging est un benchmark
très fort. Avec rebalancement discret, son erreur n'est pas nulle mais elle
décroît quand le nombre de dates augmente. Un réseau profond bien entraîné
doit obtenir une MSE comparable, parfois légèrement meilleure sur la
distribution simulée car il optimise directement l'objectif discret sous
$\mathbb{P}$. Le MLP simple peut sous-ajuster, tandis que le LSTM n'apporte pas
nécessairement de gain en modèle markovien mais peut stabiliser certains cas de
frictions.

\section{Analyse de robustesse}

La robustesse est étudiée en faisant varier :
\begin{itemize}
\item le strike $K$, qui déplace l'option entre out-of-the-money,
at-the-money et in-the-money ;
\item la volatilité $\sigma$, qui modifie la convexité effective et la fréquence
des grands mouvements ;
\item la maturité $T$, qui change l'horizon d'accumulation de l'erreur ;
\item le nombre de dates $n$, qui contrôle la granularité de rebalancement.
\end{itemize}

Les intervalles de confiance sur la MSE sont calculés à partir des carrés de
P\&L hors-échantillon :
\[
\widehat{\mathrm{MSE}}\pm1.96
\frac{\widehat{\mathrm{sd}}(\Pi_T^2)}{\sqrt{N}}.
\]
Cette approximation normale est acceptable pour de grands échantillons, mais
les queues de distribution de Put peuvent justifier un bootstrap dans une étude
plus poussée.

\section{Extensions : coûts proportionnels et CVaR}

L'extension principale implémentée est le coût proportionnel
$C_T(\delta)$. Il rend le problème plus proche de l'article original, qui met
l'accent sur les frictions de marché. La stratégie optimale ne cherche alors
plus à répliquer agressivement le delta instantané : elle arbitre réduction de
risque et turnover.

Le code inclut aussi une fonction de perte CVaR empirique :
\[
\mathrm{CVaR}_\alpha(L)=q_\alpha(L)+
\frac{1}{1-\alpha}\mathbb{E}[(L-q_\alpha(L))^+],
\qquad L=-\Pi_T.
\]
Elle n'est pas l'objectif principal demandé, mais elle permet d'étudier une
préférence plus asymétrique, orientée vers les pertes extrêmes.

\section{Conclusion}

Le projet montre comment adapter rigoureusement le Deep Hedging à un Put avec
prime connue et perte quadratique. L'approche transforme un problème de
couverture discrète en apprentissage supervisé par simulation de P\&L terminal,
où les actions sont les positions de hedge. Ses avantages sont la flexibilité
face aux coûts, contraintes et objectifs non classiques, ainsi que
l'optimisation directe hors-échantillon. Ses limites sont la dépendance à la
qualité du générateur de scénarios, le risque de surapprentissage à une
dynamique simulée et l'absence de garantie globale d'optimalité pour les
réseaux entraînés par gradient.

Dans le cadre Black--Scholes frictionless, le delta hedging reste une référence
structurellement difficile à battre. L'intérêt du Deep Hedging devient plus net
lorsque la couverture est discrète, que la perte est non standard ou que des
frictions comme les coûts proportionnels empêchent la réplication classique.

\end{document}
